\section{Task 6: The final report, the presentation and the rehearsal}
\sectionmark{Task 6: Report, presentation, rehearsal}
\label{sec:t6}

\subsection{Purpose, conditions and criterion}

\begin{plain}{Task 6 at a glance}
\textbf{Purpose.} The project is judged by what is handed in on 1 December: a report, a
presentation with a live demonstration, and the repository. This task assembles them from the
records of Tasks~1 to~5, rehearses the demonstration twice on 30 November, and marks the
version that is handed in.

\textbf{Before you begin.} The official requirements of the course for the report and the
presentation are at hand (Step~\ref{st:reqs}); this section cannot replace them, and where they
differ from it, they prevail. For the sections on results: part~B of Task~4 is complete. For
the rehearsal: Tasks~1 to~5 are complete.

\textbf{Complete when.} (1)~The report contains every item that the course requires, and every
figure in it can be traced to a record of this document. (2)~The presentation has been timed
within its limit. (3)~The demonstration has been rehearsed twice, each fallback has been
practised once, and the times are recorded. (4)~The repository carries the tag \texttt{v1.0}. (5)~Everything has
been submitted where the course requires it, and the board has been returned.

\textbf{Time.} Writing as required; each rehearsal about half an hour.
\end{plain}

\subsection{Establish the requirements}

\step{Obtain the official requirements and fill in the checklist}\label{st:reqs}
Collect the course's instructions for the final report and the final presentation from the
course site, the syllabus and the instructor's announcements, and enter them in
Table~\ref{tab:reqs}. Where an item cannot be found, ask the teaching staff now and not on the
last day.

\begin{table}[H]
\centering
\begin{tabular}{L{6.0cm}L{5.2cm}}
\toprule
\multicolumn{1}{c}{\hd{Question}} & \multicolumn{1}{c}{\hd{Answer of the course}} \\
\midrule
Report: template to be used; file format & \\
Report: required sections; limit of pages or words & \\
Report: statement of each member's contribution? & \\
Report: where and by what time it is submitted & \\
Presentation: length; number of speakers & \\
Presentation: slides required? format; where submitted & \\
Demonstration: room, time, table, power sockets & \\
Demonstration: is a projector or a second screen provided? & \\
Video: where the link is entered & \\
Repository: must the staff have access? how? & \\
Board: when and where it is returned & \\
\bottomrule
\end{tabular}
\caption{The requirements of the course, to be filled in before writing begins.}
\label{tab:reqs}
\end{table}
\checkpoint{Every row has an answer, or the name of the person who has been asked.}

\subsection{Collect the evidence}

A report of measurements is only as good as its sources. Before any text is written, the
material of Table~\ref{tab:evidence} is brought together in one shared folder.

\begin{table}[H]
\centering
\begin{tabular}{L{4.6cm}L{6.6cm}}
\toprule
\multicolumn{1}{c}{\hd{Material}} & \multicolumn{1}{c}{\hd{Where it comes from}} \\
\midrule
Geometry of the stand, scale, capital height & Record of Task~1 (Table~\ref{tab:t1record}) \\
Photograph of the stand & Task~1, taken with a telephone \\
A photograph as the board delivers it & Any file of \file{test\_set\bs images}, for example
  \file{T26.png} \\
A sheet photograph with its grid & Any file of \file{data\bs sheets\bs debug} \\
Number of photographs and crops; accuracy on real crops before and after retraining; most
  common confusions & Record of Task~2 (Table~\ref{tab:t2record}) \\
Accuracy for each class & The file \file{results\bs crops\_symbols\_v1.csv}, which opens in
  Excel \\
Size and training of the network & The end of the file \file{models\bs symbols\_v1.json}
  (open it with Notepad) \\
Kinds of error; trials of settings; decision of the gate & Record of Task~3
  (Tables~\ref{tab:triage}, \ref{tab:trialrec} and~\ref{tab:t3record}) \\
Final results against the targets & Table~\ref{tab:targets} \\
Misread lines as examples & The review document of Task~4, part~B \\
A PDF produced by the system & The file \file{out\bs document.pdf} after the recording of
  the video \\
The video & The link of Task~5 \\
The version handed in & The tag of Step~\ref{st:release} \\
\bottomrule
\end{tabular}
\caption{The evidence on which the report rests.}
\label{tab:evidence}
\end{table}

\step{Open the project in PowerShell}
% The opening sequence, written out in full. Input wherever a task may begin in a new window.
If a PowerShell window whose prompt begins with \texttt{(.venv)} is still open, use it and go
on to the next step. Otherwise open a new window (click \ui{Start}, type \code|powershell|,
click \ui{Windows PowerShell}) and give the two commands of the opening sequence in it.
\begin{ps}[Opening sequence, command 1 of 2]
cd C:\ee105\camera-to-latex
\end{ps}
\begin{ps}[Opening sequence, command 2 of 2]
.\.venv\Scripts\Activate.ps1
\end{ps}
\checkpoint{The prompt reads:}
\venvprompt


\step{Copy the sample output into the repository}
The PDF that the system produced during the recording of the video is the most direct proof
that it works. Copy it into the folder of documents and record it:
\begin{ps}
Copy-Item out\document.pdf docs\sample_output.pdf
\end{ps}
\begin{ps}
git add docs\sample_output.pdf
\end{ps}
\begin{ps}
git commit -m "Sample output produced from photographs"
\end{ps}
\begin{ps}
git push
\end{ps}
\checkpoint{The file \file{sample\_output.pdf} is visible in the folder \file{docs} on the
GitHub page of the repository.}

\subsection{Write the report}

If the course provides a template, its sections are used. Table~\ref{tab:outline} lists what a
complete account of this project contains; each item belongs somewhere in the template. It
follows the outline of the project plan.

\begin{longtable}{C{0.6cm}L{3.0cm}L{7.3cm}}
\toprule
 & \multicolumn{1}{c}{\hd{Part}} & \multicolumn{1}{c}{\hd{What it must say}} \\
\midrule
\endhead
\bottomrule
\endfoot
\bottomrule
\caption{Contents of a complete report.}\label{tab:outline}\\
\endlastfoot
1 & Problem and scope & What the system does (one printed line per photograph, read into
  \LaTeX{} and collected into a PDF) and what it does not do (handwriting, several lines at
  once, recognition on the board itself) \\
2 & Hardware and its limits & The board and the camera; 256\,KB of memory, hence photographs of
  $320 \times 240$ gray pixels; the stand; the measured scale and the size of the letters \\
3 & Architecture & The stages from capture to PDF, with the diagram of the pipeline; why
  recognition runs on the laptop \\
4 & Data & The synthetic images; the 54 sheet photographs and their crops; the 60 test and 20
  development lines, and the rule that kept them apart \\
5 & Classifier & The network (159 classes for 182 symbols, 149{,}311 weights); its accuracy on
  held-out real crops before and after retraining, as counts; the common confusions \\
6 & Cutting and layout & How ink is separated from paper and formed into symbols; how
  superscripts, subscripts and fractions follow from positions and sizes \\
7 & Results & Table~\ref{tab:targets}: every target, the result as a count, whether it was
  met, and the precision of the figure \\
8 & Failures & Three or four misread lines, each with its photograph, the kind of error and the
  reason \\
9 & Limits and next steps & What a second version would change, for example recognition on the
  board, or a model that reads a whole line \\
10 & Contributions & Who did what, in the form that the course requires \\
11 & Sources & The software libraries and fonts used, and any assistance, including that of AI
  tools, in the form that the course requires \\
\end{longtable}

\step{Write the parts that do not depend on the final results}
Parts 1 to 4, 6 and 9 to 11 can be written before part~B of Task~4 is finished.
\checkpoint{These parts exist in draft, and every number in them is taken from a record or
from the repository, not from memory.}

\step{Write the results as they were measured}
Parts 5, 7 and 8 are written from Table~\ref{tab:targets} and the review document. Four rules
keep them honest and useful.
\begin{enumerate}
\item A result is given as a count and a percentage, for example ``17 of 25 lines (68\,\%)'',
  and a target that was missed is called missed.
\item Simulated and real photographs are never mixed in one figure, and each figure says which
  it is. The results of 23 November were obtained on simulated photographs.
\item The uncertainty is named where a conclusion depends on it; the formula and the values
  are in Section~\ref{sec:stat}.
\item The conditions are stated: print of 40\,pt (or 48\,pt, if the larger print was used), the
  measured distance and scale, the lighting, and the fact that each line was read in the mode,
  text or mathematics, that the operator would choose.
\end{enumerate}
\checkpoint{Every figure of the results can be found in Table~\ref{tab:targets} or in a record
of this document.}

\step{Check the report against the requirements}
Read the finished report beside Table~\ref{tab:reqs} and Table~\ref{tab:outline}.
\checkpoint{Every required section is present, the limit of length is kept, the link of the
video opens in a private browser window, and a member who did not write the report has read it
from beginning to end.}

\subsection{Prepare the presentation}

\step{Bring the slides up to date}
The team has a deck of slides from the design presentation, \emph{Camera to \LaTeX{}: how it
works and why}; the project lead shares it. Its account of the method remains valid. Three
things change for the final presentation: figures from simulation are replaced by the measured
ones; the stand and the results are shown; and time is reserved for the demonstration.
Table~\ref{tab:slides} proposes a sequence, to be shortened to the length that the course
allows.

\begin{table}[H]
\centering
\begin{tabular}{C{1.1cm}L{4.2cm}L{5.7cm}}
\toprule
\hd{Slide} & \multicolumn{1}{c}{\hd{Subject}} & \multicolumn{1}{c}{\hd{Content}} \\
\midrule
1 & Title & Project, team, date \\
2 & The task and the limits & One sentence on the aim; the three limits that shaped the design:
  memory, resolution, one line per photograph \\
3 & How it works & The diagram of the pipeline \\
4 & Data and classifier & Synthetic images and sheet photographs; accuracy on real crops,
  before and after \\
5 & The stand & Its photograph; distance, scale, size of the letters \\
6 & Results & Table~\ref{tab:targets}, as counts, with the targets \\
7 & What fails, and why & Two misread lines with their photographs \\
8 & Demonstration & Live; see Section~\ref{sec:rehearsal} \\
9 & Limits and next steps & Two or three points \\
\bottomrule
\end{tabular}
\caption{A sequence of slides for the final presentation.}
\label{tab:slides}
\end{table}
\checkpoint{The slides contain no figure from simulation that is presented as a result, and a
copy of the slides as a PDF is stored on the demonstration laptop, so that the presentation
does not depend on a network connection.}

\step{Time the presentation}
Speak the whole presentation aloud once, with a stopwatch, without the demonstration.
\checkpoint{The spoken part and the demonstration together fit into the time that the course
allows, with half a minute to spare.}

\subsection{Rehearse the demonstration}
\label{sec:rehearsal}

The demonstration is rehearsed twice on 30 November, one of the two times, if possible, in the
room
of the presentation, where the light and the table differ from those at home.

\step{Pack}
\begin{table}[H]
\centering
\begin{tabular}{L{5.5cm}L{5.7cm}}
\toprule
\multicolumn{1}{c}{\hd{Item}} & \multicolumn{1}{c}{\hd{Remark}} \\
\midrule
\cbox\ \ Board in its box & \\
\cbox\ \ USB data cable, and a second one & The cable that passed the soak test, and a spare \\
\cbox\ \ Laptop and charger & The laptop of Tasks~1 to~4 \\
\cbox\ \ Books of the stand, the arm, two cloths & Or the list of their heights from the record
  of Task~1 \\
\cbox\ \ Line platform & Marked in Task~1 \\
\cbox\ \ Lamps, extension lead & The room's light is unknown \\
\cbox\ \ Blank white sheets, tape, ruler & \\
\cbox\ \ The three cards, and second copies & Table~\ref{tab:cards} \\
\cbox\ \ Record of Task~1 & Height of the lens, thickness of the platform \\
\cbox\ \ Slides as a PDF, and the video as a file & Both on the laptop, usable without a
  network \\
\bottomrule
\end{tabular}
\caption{What is taken to the rehearsal and to the presentation.}
\label{tab:pack}
\end{table}

\step{Set up against the clock}\label{st:setup6}
Start a stopwatch and build the system as it will be built on the day.
\begin{enumerate}
\item Build the stand as in Task~1. Check with the ruler that the lens is at the height of the
  record, and put the line platform in place with a white sheet on it.
\item Place the lamps, connect the board, open a PowerShell window, and give the opening
  sequence:
\begin{ps}[Opening sequence, command 1 of 2]
cd C:\ee105\camera-to-latex
\end{ps}
\begin{ps}[Opening sequence, command 2 of 2]
.\.venv\Scripts\Activate.ps1
\end{ps}
\item Confirm that the board answers:
\begin{ps}
python -m pipeline.capture --ping
\end{ps}
\item Check the picture with the widest line. Lay page~8 below the camera, centre line T39,
  look for shadows, and close the viewfinder with \key{Q}:
\begin{ps}
python -m pipeline.capture --series data\raw\rig --names "room_check"
\end{ps}
\item Start the application:
\begin{ps}
python -m pipeline.main --new
\end{ps}
\end{enumerate}
\checkpoint{The board answered, line T39 was wholly inside the picture and free of shadows, and
the application shows the live picture. Stop the stopwatch and enter the time in the record;
the aim is ten minutes or less.}

\step{Run the demonstration}\label{st:demo}
The live demonstration follows the sequence of the video (Table~\ref{tab:storyboard}) and adds
two things that a video of one minute cannot show: a correction by hand, and the removal of a
line.
\begin{enumerate}
\item Read and accept the three cards as in the video: \key{M} for the mode, \key{Space},
  \key{A}.
\item \textbf{Correction.} Read one further line, of your choice, and press \key{E} in place of
  \key{A}. The PowerShell window prints the reading after the word \emph{current} and waits.
  Click on the PowerShell window, type the corrected \LaTeX{} text (or press \key{Enter} alone
  to keep the reading), and press \key{Enter}. The line is accepted if it compiles. Then click
  on the picture of the application window again, so that the keys act on it. While PowerShell
  waits for the text, the application window does not respond; that is expected.
\item \textbf{Removal.} Press \key{U}. PowerShell prints \emph{removed}, followed by the line
  that was accepted last.
\item Press \key{P}. The PDF opens with the three lines.
\end{enumerate}
\checkpoint{The PDF shows the three lines. Enter the duration of the demonstration in the
record; the aim is three minutes or less.}

\begin{table}[H]
\centering
\begin{tabular}{C{2.2cm}L{8.8cm}}
\toprule
\hd{Key} & \multicolumn{1}{c}{\hd{Effect in the application window}} \\
\midrule
\key{Space} & Read the line that is in view; the picture freezes and shows boxes and reading \\
\key{A} or \key{Enter} & Accept the reading into the document, if it compiles \\
\key{E} & Edit the reading in the PowerShell window, then accept it \\
\key{R} & Reject the reading and return to the live picture \\
\key{M} & Change the mode: automatic, text, mathematics \\
\key{U} & Remove the line that was accepted last \\
\key{P} & Build the PDF and open it \\
\key{S} & Save the current picture in the folder \file{data\bs raw} \\
\key{Q} or \key{Esc} & Leave the application; the PDF is built if lines were accepted \\
\bottomrule
\end{tabular}
\caption{Keys of the application. They are lower-case letters and act only while the
application window is the active one.}
\label{tab:appkeys}
\end{table}

\step{Practise the failures}\label{st:drills}
A demonstration fails gracefully only if its failures have been practised. Carry out each of
the following once in the first rehearsal.
\begin{enumerate}
\item \textbf{A wrong reading.} Press \key{R}, move the card slightly, and press \key{Space}
  again. If the second reading is wrong too, correct it with \key{E} and say what was done;
  this shows a feature and is no failure.
\item \textbf{The board does not answer.} Leave the application with \key{Q}, unplug the cable
  and plug it in again, repeat the ping command and the application command of
  Step~\ref{st:setup6}.
\item \textbf{The board cannot be revived: first fallback.} Read the stored photographs of the
  three cards with the three commands of Section~\ref{sec:noboard}, and say that stored
  photographs are being read.
\item \textbf{The laptop fails: second fallback.} Play the video, from the copy on a second
  device.
\end{enumerate}
\checkpoint{Each of the four has been carried out once, and each member knows who does what in
each case.}

\step{Rehearse a second time, and take down the stand in order}
Repeat Steps~\ref{st:setup6} and~\ref{st:demo} from the packed state. Afterwards write down
anything that had to be improvised, pack the items of Table~\ref{tab:pack}, and store the board
in its box.
\checkpoint{The record contains the times of both rehearsals.}

\subsection{Mark and hand in the final version}

\step{Mark the version that is handed in}\label{st:release}
After the second rehearsal, when nothing further will change:
\begin{ps}
git pull
\end{ps}
\begin{ps}
python -m pytest -q
\end{ps}
\begin{ps}
git status --short
\end{ps}
\checkpoint{The tests pass, and the third command prints nothing.}
\begin{ps}
git tag -a v1.0 -m "Version handed in on 1 December 2026"
\end{ps}
\begin{ps}
git push origin v1.0
\end{ps}
\checkpoint{The GitHub page of the repository lists the tag \texttt{v1.0} under \emph{Tags}.}

\step{Give the staff access to the repository, if required}
The repository is private. If Table~\ref{tab:reqs} says that the staff must see it, the lead
either invites them (on the GitHub page: \menu{Settings}\then\menu{Collaborators}\then\menu{Add
people}, with their GitHub user names) or hands in a copy: on the page of the repository,
choose the tag \texttt{v1.0} in the branch menu, then \menu{Code}\then\menu{Download ZIP}.
\checkpoint{The staff have confirmed access, or the ZIP file has been submitted.}

\step{Submit, and go through the final checklist}
\begin{table}[H]
\centering
\begin{tabular}{L{11.2cm}}
\toprule
\multicolumn{1}{c}{\hd{Before the deadline of 1 December}} \\
\midrule
\cbox\ \ Report submitted in the required format at the required place \\
\cbox\ \ Slides submitted, if required, and stored on the demonstration laptop \\
\cbox\ \ Link of the video entered where required, and opened once more in a private window \\
\cbox\ \ Tag \texttt{v1.0} uploaded; access or ZIP provided, if required \\
\cbox\ \ Table of results in the README up to date \\
\cbox\ \ Items of Table~\ref{tab:pack} packed for the presentation \\
\cbox\ \ Board returned as the course requires, after the presentation \\
\bottomrule
\end{tabular}
\caption{Final checklist.}
\label{tab:final}
\end{table}

\step{Fill in the record}
\begin{table}[H]
\centering
\begin{tabular}{L{6.4cm}C{2.3cm}C{2.3cm}}
\toprule
\multicolumn{1}{c}{\hd{Quantity}} & \hd{Rehearsal 1} & \hd{Rehearsal 2} \\
\midrule
Place; date and time & & \\
Time to set up & \hfill min & \hfill min \\
Duration of the demonstration & \hfill min & \hfill min \\
Cards read correctly at the first attempt (of 3) & & \\
Failures practised (of 4) & & -- \\
Duration of the spoken presentation & \hfill min & \hfill min \\
\bottomrule
\end{tabular}
\caption{Record of Task 6.}
\label{tab:t6record}
\end{table}
